\subsection{Further star reductions at eighteen arrows}

\begin{lemma}[Proved: merging a trivial leaf into an arbitrary leaf]
\label{CP2:A-star-general-leaf-merge}
Consider a source star with center $M$, one leaf $N$ on a
left-$N$, right-$M$ biset $S$, and a trivial leaf on a right
$M$-act $T$. If there is an $M$-equivariant map $T\to S$,
the star retracts onto its full $(M,N,S)$ edge. If the star
has the rational homology and cohomology ring of $\CP^2$,
so does that edge.
\end{lemma}
\begin{proof}
Map the trivial leaf to the object with corner $N$, its
identity to $1_N$, and its stalk by the specified map.
Keep the retained edge unchanged. Equivariance is exactly
the remaining functoriality condition, so this is an actual
category retraction.

In the double mapping-cylinder model its quotient by the
retained edge is the unreduced suspension of $B(T,M)$:
the center, the retained leaf and its entire bridge are
collapsed, leaving the other bridge with its two ends
collapsed separately. If the retract had zero second
cohomology, every second class on the star would lift from
this suspension and all second products would be zero.
Thus its second group is nonzero. The retraction makes its
homology a direct summand of the rational projective-plane
homology. Pullback in degree two is consequently an
isomorphism; the nonzero square forces its fourth group
to be nonzero as well. These facts give exactly the stated
rational homology and ring. No condition on $BN$ is used.
\end{proof}

\begin{lemma}[Proved: a constant actor on a connected leaf action space]
\label{CP2:A-connected-leaf-constant}
Let $S$ be a left-$N$, right-$M$ biset, where $BN$ is
rationally acyclic. Suppose the ordinary left action space
$Z=B(N,S)$ is connected and some right $M$ actor on $S$
is constant. The full two-object collage $(M,N,S)$ has
zero rational second homology.
\end{lemma}
\begin{proof}
Collapsing $BN$ in the double-bar pushout is a rational
homology equivalence onto
$Y=\operatorname{hocofib}(Z_{hM}\to BM)$.
The cone of the augmented right $\mathbb QM$ chain complex
$C_*(Z;\mathbb Q)\to\mathbb Q$ has zero homology in
degrees zero and one, since $Z$ is nonempty and connected.
Its degree-two homology is $H_1(Z;\mathbb Q)$.
Tensoring a projective left monoid resolution gives the
relative bar complex for $Y$. Its lowest-degree homology,
or equivalently the first nonzero row of the hyperhomology
spectral sequence, therefore gives
\[
 H_2(Y;\mathbb Q)=H_1(Z;\mathbb Q)_M.
\]
If $m$ acts constantly with value $u$, commutation makes
$u$ fixed by every $N$ actor. The induced map on $Z$
factors as the forgetful map $Z\to BN$ followed by the
fixed-state inclusion $BN\to Z$. It acts as zero on
$H_1(Z;\mathbb Q)$. The coinvariant relation $v m=v$
then kills every class $v$, proving the claim.
\end{proof}

\begin{proposition}[Proved: a binary leaf beside a trivial leaf]
\label{CP2:A-star-binary-leaf-eleven}
A source star with center of order at most eleven, a
trivial leaf on any nonempty finite stalk, and another
leaf on at most two states cannot have the rational
homology and cohomology ring of $\CP^2$. There is no
restriction on the other leaf corner.
\end{proposition}
\begin{proof}
Write $S$ for the binary stalk and $N$ for its leaf corner.
A singleton stalk has a center-fixed state; more generally,
if the left $N$ image contains a constant map, its value is
fixed by the commuting center action. The constant map
from the other stalk to that point is center-equivariant.
Lemma~\ref{CP2:A-star-general-leaf-merge} transfers the
rational projective-plane target to the retained binary
edge, contradicting the universal two-cross obstruction.

If the left image contains a transposition but no constant,
both corner actions on $S$ are permutations. Its bridge
maps as a finite covering to $BN\times BM$. Transfer
makes rational cohomology pullback injective. Compatibility
of the restrictions of a global second class to the two
corners, together with $H^1(BN;\mathbb Q)=H^1(BM;\mathbb Q)=0$,
therefore forces both restrictions to be zero. The other
leaf is trivial, so the two-block relative-cup lemma kills
all second products.

The remaining left image is the identity. Apply
Lemma~\ref{CP2:A-trivial-cut-H2} to replace $N$ by a
trivial corner. This is an actual retract and an $H^2$
isomorphism, so its homology is bounded. Merging its two
trivial leaves gives the ordinary center-act cone on the
disjoint union of the stalks, homotopy equivalent to the
retracted star wedged with a circle. The bounded-eleven
act theorem kills every second product there. Pullback
along the $H^2$ isomorphism gives the contradiction.
\end{proof}

\begin{proposition}[Proved: a small three-state leaf]
\label{CP2:A-star-small-three-leaf}
A source star with center of order at most eleven, one
trivial leaf, and another leaf of order at most four on
at most three states cannot have the rational homology
and cohomology ring of $\CP^2$. The other stalk may have
any positive size. The same conclusion holds if the
otherwise trivial leaf has order at most two.
\end{proposition}
\begin{proof}
At a leaf of order at most two, exact orbit comparison
replaces it by a trivial leaf without increasing its
stalk. We may therefore assume that leaf trivial.
The binary case is the preceding proposition. Write
$N$ for the small leaf corner and $S$ for its three-set.
If its trivial right rational module is projective,
exact orbit comparison makes this leaf trivial as well;
merging the leaves and applying the bounded-eleven theorem
gives the result.

Otherwise the small minimal-ideal classification makes
the minimal ideal of $N$ a one-row band with at least
two columns and trivial maximal group. Its left actors
on $S$ are retractions onto a common nonempty set $U$,
fixed pointwise by every actor of $N$. In particular
$BN$ is contractible. If $|U|=1$, its point is fixed by
the center. Restrictions of every global second class
to the center are zero, as are those to both leaves;
the two-block lemma applies. If $|U|=3$, the entire
$N$ action is trivial, so the leaf can be collapsed by
a homotopy equivalence and the bounded-eleven argument
applies.

Let $|U|=2$. If all band actors are the same retraction
$P$, then $P$ commutes with every actor of $N$: left
multiplication fixes each band element, and right
multiplication remains in the band. It also commutes
with the center. Retracting the stalk to $U$ transfers
the rational target by Lemma~\ref{CP2:A-star-leaf-reductions}(4).
On $U$ the leaf action is trivial. Collapse its contractible
corner, merge the now trivial leaves, and use bounded-eleven.

Finally suppose two band actors $P,Q$ differ. For the
unique state $w\notin U$, their values are the two distinct
points of $U$. A transformation commuting with $P,Q$ is
either the identity on $S$ or a constant map to one of
these two points: the image of $w$ determines all three
values. If the center image consists only of the identity,
there is a center-fixed state and the relative-cup argument
already applies. Otherwise a center actor is constant.
The left $N$ action category is connected, since its arrows
from $w$ reach both points of $U$. The preceding connected
leaf lemma makes the full $(M,N,S)$ edge have $H^2=0$.
Its union with the remaining trivial leaf is then a
two-block collage with both second restrictions zero.
This finishes all cases.
\end{proof}

\begin{proposition}[Proved: four states opposite an order-three leaf]
\label{CP2:A-star-three-leaf-four-states}
A source star with center of order at most eight, one
trivial leaf, and another leaf of order at most three on
at most four states cannot have the rational homology
and cohomology ring of $\CP^2$.
\end{proposition}
\begin{proof}
The preceding proposition treats at most three states.
An order-three leaf with projective trivial right module
is replaced by its exact orbit set, and the merged
ordinary act is excluded by bounded-eleven. Otherwise
the leaf is $R_3=\{1,p,q\}$, with $ab=b$ on its two
nonidentity elements. Its actors $P,Q$ are retractions
onto the same subset $U$ of the four-state stalk.

If $|U|=1$, a center-fixed state kills the products. If
$P=Q$, their common biset retraction transfers the rational
target to the stalk $U$. The leaf now acts trivially;
collapse its contractible corner and use bounded-eleven.
If $|U|=3$ and $P\ne Q$, the natural graph model of
Lemma~\ref{CP2:A-high-P5-graph} has a single connecting
edge and one isolated vertex. Equivariant collapse to
its two contractible components replaces the entire
leaf and stalk by a trivial two-state leaf, by homotopy
pushout invariance. Again merge and use bounded-eleven.

It remains that $|U|=2$, $P\ne Q$. The natural graph
has two vertices and two edges; it is connected and
has one cycle. By Lemma~\ref{CP2:A-one-cycle-graph-cofiber},
if any center actor has cycle degree different from one,
the retained full edge has zero $H^2$, and its union
with the trivial leaf has zero second products.
Assume every actor has degree one.

If the graph has one loop and one connecting edge, all
actors fix the loop vertex. If its two connecting edges
have the same direction, degree one forces both vertices
fixed. Either case gives a center-fixed state. In the
last case the two edges have opposite directions and
the only possible actors are the identity and half-turn,
as in the four-state graph proof cited above.

If the center has $H^2(BM;\mathbb Q)=0$, both sides of
the cut from the center to its two leaves have zero
second restriction, and the two-block lemma applies.
Otherwise the center's minimal ideal has at least two
rows and two columns. A nontrivial maximal group would
require at least eight ideal elements and an additional
identity, contrary to $|M|\leq8$. Its maximal group is
therefore trivial. That local group surjects onto the
universal group of $M$, so every homomorphism from $M$
to the half-turn group is trivial. All actors fix both
vertices, giving the remaining center-fixed state.
This proves the assertion.
\end{proof}

\begin{proposition}[Proved: a twelve-element center with two binary leaves]
\label{CP2:A-star-twelve-binary-leaves}
A source star with a center of order twelve, two trivial
leaves and two arrows to each leaf cannot have type
$\CP^2$; indeed it cannot have its rational homology and ring.
\end{proposition}
\begin{proof}
Merge the two leaves. The resulting ordinary act cone
on the disjoint union of two invariant two-sets has
bounded rational homology and the same products in
degrees at least two. A possible nonzero product requires
at least three minimal rows and two columns, by the
two-row act obstruction; its maximal group is trivial
by the order budget. The possible minimal bands are
$3\times2$, $4\times2$, $5\times2$, and $3\times3$,
with respectively six, four, two, and three outside
elements. The first three are excluded by
Propositions~\ref{CP2:A-three-row-six-outside-acts} and
\ref{CP2:A-two-column-four-outside-bounded}.

For the $3\times3$ band, a band-fixed state is fixed by
the whole monoid and kills products. In its absence
each invariant two-set is a full core orbit. Their
column maps give two two-part partitions of the three
band columns. If the partitions differ, the constant
action quotient identifies all four states, since
the bipartite graph on the four parts is connected.
Its cone has zero second products; the natural $H^2$
surjection from the constant-action quotient kills
those on the original cone.

If the partitions agree, collapsing columns by their
common two-part partition is a monoid congruence:
every outside column actor preserves that partition,
as witnessed on either core orbit. The entire four-state
action factors through the quotient monoid, which has
minimal band $3\times2$ and three outside elements.
Lemma~\ref{CP2:A-column-quotient-act-H2} gives an $H^2$
isomorphism. The three-row act theorem kills products
on the quotient without any boundedness assumption,
and naturality finishes the proof.
\end{proof}

\begin{remark*}[Exact scope of the new large-corner reductions]
In the notation $(a,b,c,s,t)$ for center, leaf orders and
the corresponding stalk sizes, these propositions exclude
the following seven eighteen-arrow profiles and their
leaf interchanges:
\[
\begin{gathered}
 (8,1,3,2,4),\quad(8,1,3,3,3),\quad(8,1,4,2,3),\\
 (8,1,5,2,2),\quad(8,2,3,2,3),\quad(9,1,3,2,3),\\
 (12,1,1,2,2).
\end{gathered}
\]
The two coned-three-state profiles $(5,1,8,1,3)$ and
$(8,1,5,1,3)$ remain separate obligations here.
\end{remark*}
